% ============================================================
% semana_03_conjuntos — teoria
% ============================================================
% Este archivo es incluido desde main.tex con \include{}


% --- TEORÍA: TEORÍA DE CONJUNTOS -------------------------------------------
%  Curso: Aritmética | Semana: 01 | Unidad 1: Fundamentos Numéricos

% --- PORTADA DE CLASE ------------------------------------------------------
\portadaclase{Teor\'{\i}a de Conjuntos}{01}{1}{Fundamentos Num\'ericos}

% --- MOTIVACIÓN ------------------------------------------------------------
\section*{¿Por qu\'e estudiar este tema?}
\begin{tcolorbox}[
  enhanced, colback=AzulClaro!30, colframe=AzulMedio,
  arc=3pt, boxrule=0.6pt, left=10pt, right=10pt, top=6pt, bottom=6pt
]
\textit{La Teor\'{\i}a de Conjuntos es la base estructural de toda la matem\'atica
moderna. En los ex\'amenes de admisi\'on de la UNMSM, UNI, UNSCH, PUCP y Villarreal
aparece de forma directa en Aritm\'etica, \'Algebra y Razonamiento Matem\'atico.
Dominar operaciones entre conjuntos y conjuntos potencia es indispensable para
resolver hasta 6 preguntas en un solo examen.}
\end{tcolorbox}
\vspace{0.5cm}

% --- SECCIÓN 1: NOCIÓN DE CONJUNTO -----------------------------------------
\section{Noci\'on de Conjunto}

\subsection{Definici\'on}

\begin{cajadef}
  \textbf{Conjunto:} Reunión, colección o agrupación de objetos abstractos y/o
  concretos, que pueden o no tener una caracter\'{\i}stica en com\'un. A cada objeto
  perteneciente al conjunto se le llama \textbf{elemento}.\\[4pt]
  \textit{Ejemplos:} el conjunto de las vocales, el conjunto de los n\'umeros
  naturales menores que 10, el conjunto de los colores primarios.
\end{cajadef}

\subsection{Determinaci\'on de un conjunto}

\begin{cajaprop}[A --- Por comprensi\'on (constructiva)]
  Se describe una caracter\'{\i}stica com\'un que identifica a todos los elementos.
  \[
    A = \{\,x \mid x \text{ es una vocal}\,\}
  \]
\end{cajaprop}

\begin{cajaprop}[B --- Por extensi\'on (tabular)]
  Se listan uno a uno todos los elementos del conjunto.
  \[
    A = \{a,\; e,\; i,\; o,\; u\}
  \]
\end{cajaprop}

\begin{cajatip}
  En los exámenes, la notación por comprensión usa la barra vertical ``$\mid$''
  o la diagonal ``/''. Ambas significan ``tal que''. Así,
  $\{x \mid P(x)\} = \{x / P(x)\}$.
\end{cajatip}

% --- SECCIÓN 2: RELACIÓN DE PERTENENCIA ------------------------------------
\section{Relaci\'on de Pertenencia}

\begin{cajadef}
  \textbf{Pertenencia ($\in$):} Un objeto $a$ pertenece al conjunto $A$ si $a$
  es un elemento de $A$. Se escribe $a \in A$. En caso contrario, $a \notin A$.
\end{cajadef}

\begin{cajaejemplo}[1]
  Dado $A = \{1,\; 3,\; \{4\},\; 8\}$, analizar:
  \begin{align*}
    1 &\in A && \checkmark \quad \text{(1 es elemento directo de } A\text{)} \\
    \{4\} &\in A && \checkmark \quad \text{(\{4\} aparece listado en } A\text{)} \\
    9 &\notin A && \checkmark \quad \text{(9 no está en la lista)} \\
    \{8\} &\notin A && \checkmark \quad \text{(8 sí pertenece, pero \{8\} como conjunto, no)}
  \end{align*}
\end{cajaejemplo}

\begin{cajaadvert}
  \begin{itemize}
    \item \textbf{Error frecuente:} confundir el elemento $8$ con el conjunto
      $\{8\}$. Son objetos distintos: $8 \in A$ pero $\{8\} \notin A$ para el
      conjunto del ejemplo anterior.
    \item \textbf{Error frecuente:} escribir $\{a, b\} \in C$ cuando en realidad
      $\{a, b\}$ es un subconjunto. Verificar si ese par aparece \emph{listado
      directamente} como elemento.
  \end{itemize}
\end{cajaadvert}

% --- SECCIÓN 3: CARDINAL ---------------------------------------------------
\section{Cardinal de un Conjunto}

\begin{cajadef}
  El \textbf{cardinal} de un conjunto $A$, denotado $n(A)$, es la cantidad de
  elementos \emph{diferentes} que contiene. Si un elemento se repite, se cuenta
  una sola vez.
\end{cajadef}

\begin{cajaejemplo}[2]
  \[
    B = \{1,\; 2,\; 2,\; 3,\; 3,\; 3\}
  \]
  Los elementos distintos son $1$, $2$ y $3$, por tanto $n(B) = 3$.
\end{cajaejemplo}

% --- SECCIÓN 4: RELACIONES ENTRE CONJUNTOS ---------------------------------
\section{Relaciones entre Conjuntos}

% \texorpdfstring: los símbolos math de unicode-math no pueden ir directo
% a los marcadores PDF de hyperref (error "Improper alphabetic constant").
\subsection{Inclusi\'on (\texorpdfstring{$\subset$}{⊂})}

\begin{cajadef}
  $A \subset B$ (``$A$ está incluido en $B$'') si \emph{todo} elemento de $A$
  también es elemento de $B$. Es decir:
  \[
    A \subset B \;\Longleftrightarrow\; \forall\, x:\; x \in A \Rightarrow x \in B
  \]
  Si existe al menos un elemento de $B$ que no pertenece a $A$, entonces
  $A \subsetneq B$ (inclusión propia).
\end{cajadef}

\subsection{Igualdad}

\begin{cajadef}
  Dos conjuntos $A$ y $B$ son \textbf{iguales} ($A = B$) si tienen exactamente
  los mismos elementos, independientemente del orden en que se listen.\\[4pt]
  \textit{Ejemplo:} $A = \{R, O, M, A\}$ y $B = \{M, O, R, A\}$ $\Rightarrow$
  $A = B$.
\end{cajadef}

% --- SECCIÓN 5: PRINCIPALES CONJUNTOS --------------------------------------
\section{Principales Conjuntos}

\begin{cajaprop}[Conjunto vac\'{\i}o]
  No contiene ning\'un elemento. Se denota $\emptyset$ o $\{\}$.
  Su cardinal es $n(\emptyset) = 0$.\\
  \textbf{Propiedad clave:} $\emptyset \subset A$ para \emph{todo} conjunto $A$.
\end{cajaprop}

\begin{cajaprop}[Conjunto unitario o singleton]
  Posee exactamente un elemento. Ejemplo: $\{5\}$, $\{\emptyset\}$.
\end{cajaprop}

\begin{cajaprop}[Conjunto universal ($U$)]
  Es el conjunto referencial que contiene a todos los conjuntos considerados
  en un contexto dado.
\end{cajaprop}

\begin{cajaprop}[Conjunto potencia ($\mathcal{P}(A)$)]
  El conjunto de \emph{todos los subconjuntos} de $A$, incluyendo $\emptyset$ y $A$.
  \[
    n\bigl[\mathcal{P}(A)\bigr] = 2^{\,n(A)}
  \]
  \textit{Ejemplo:} Si $A = \{x, y\}$, entonces
  $\mathcal{P}(A) = \{\emptyset,\; \{x\},\; \{y\},\; \{x,y\}\}$
  y $n[\mathcal{P}(A)] = 2^{2} = 4$.
\end{cajaprop}

\subsection{Conjuntos num\'ericos}

\begin{cajaformula}
  \begin{center}
  \renewcommand{\arraystretch}{1.5}
  \begin{tabular}{@{}cl@{}}
    \toprule
    \rowcolor{AzulClaro!60}
    \textbf{S\'{\i}mbolo} & \textbf{Nombre y descripci\'on} \\
    \midrule
    $\mathbb{N}$ & Naturales: $\{0, 1, 2, 3, \ldots\}$ \\
    $\mathbb{Z}$ & Enteros: $\{\ldots, -2, -1, 0, 1, 2, \ldots\}$ \\
    $\mathbb{Q}$ & Racionales: n\'umeros expresables como $p/q$, $q \neq 0$ \\
    $\mathbb{Q}'$& Irracionales: $\sqrt{2},\; \pi,\; e,\; \ldots$ \\
    $\mathbb{R}$ & Reales: $\mathbb{Q} \cup \mathbb{Q}'$ \\
    $\mathbb{C}$ & Complejos: $a + bi$, $\; b \in \mathbb{R}$ \\
    \bottomrule
  \end{tabular}
  \end{center}
  \[
    \mathbb{N} \subset \mathbb{Z} \subset \mathbb{Q} \subset \mathbb{R}
    \subset \mathbb{C}
  \]
\end{cajaformula}

% --- SECCIÓN 6: CUANTIFICADORES --------------------------------------------
\section{Cuantificadores L\'ogicos}

\begin{cajaformula}
  \begin{center}
  \renewcommand{\arraystretch}{1.4}
  \begin{tabular}{@{}cll@{}}
    \toprule
    \rowcolor{AzulClaro!60}
    \textbf{S\'{\i}mbolo} & \textbf{Clase} & \textbf{Se lee} \\
    \midrule
    $\forall$ & Universal & ``Para todo'' \\
    $\exists$ & Existencial & ``Existe al menos un'' \\
    \bottomrule
  \end{tabular}
  \end{center}
\end{cajaformula}

\begin{cajaejemplo}[3 --- Cuantificadores con $U = \{1, 2, 3\}$]
  Determinar el valor de verdad de:
  \[
    \exists\, x,\; \forall\, y \;/\; x^{2} < y + 1
  \]
  \noindent\textit{Se pide:} que \emph{exista} un $x$ tal que para
  \emph{todo} $y$ se cumpla $x^{2} < y + 1$.
  \begin{align*}
    x = 1: &\quad 1 < 1+1=2 \;\checkmark,\quad 1 < 2+1=3\;\checkmark,\quad
              1 < 3+1=4\;\checkmark
  \end{align*}
  Como $x = 1$ satisface la condición para todo $y \in U$, la proposición
  es \textbf{Verdadera}.
\end{cajaejemplo}

% --- SECCIÓN 7: OPERACIONES ENTRE CONJUNTOS --------------------------------
\section{Operaciones entre Conjuntos}

\subsection{Uni\'on}

\begin{cajadef}
  $A \cup B = \{x \mid x \in A \;\text{o}\; x \in B\}$. Contiene todos los
  elementos que pertenecen a $A$, a $B$ o a ambos.
  \[
    n(A \cup B) = n(A) + n(B) - n(A \cap B)
  \]
\end{cajadef}

\subsection{Intersecci\'on}

\begin{cajadef}
  $A \cap B = \{x \mid x \in A \;\text{y}\; x \in B\}$. Contiene solo los
  elementos comunes a ambos conjuntos.\\
  Si $A \cap B = \emptyset$, los conjuntos se llaman \textbf{disjuntos}.
\end{cajadef}

\subsection{Diferencia}

\begin{cajadef}
  $A - B = \{x \mid x \in A \;\text{y}\; x \notin B\}$. Elementos de $A$ que
  no pertenecen a $B$.
  \[
    n(A - B) = n(A) - n(A \cap B)
  \]
\end{cajadef}

\subsection{Diferencia sim\'etrica}

\begin{cajadef}
  $A \mathbin{\triangle} B = (A - B) \cup (B - A)$. Elementos que pertenecen
  a uno de los dos conjuntos pero no a ambos.
  \[
    n(A \mathbin{\triangle} B) = n(A) + n(B) - 2\,n(A \cap B)
  \]
\end{cajadef}

\subsection{Complemento}

\begin{cajadef}
  $A' = A^{C} = U - A = \{x \in U \mid x \notin A\}$.
  \[
    n(A') = n(U) - n(A)
  \]
\end{cajadef}

\begin{cajaformula}
  \textbf{Fórmula del cardinal de la unión (tres conjuntos):}
  \begin{align*}
    n(A \cup B \cup C) =\;& n(A) + n(B) + n(C) \\
                          &- n(A \cap B) - n(A \cap C) - n(B \cap C) \\
                          &+ n(A \cap B \cap C)
  \end{align*}
\end{cajaformula}

% --- SECCIÓN 8: EJEMPLOS RESUELTOS DE NIVEL ADMISIÓN -----------------------
\section{Ejemplos resueltos de nivel admisi\'on}

\begin{cajaejemplo}[4 --- UNMSM: Conjunto potencia y cardinales]
  Se tienen dos conjuntos $A \subset B$. La diferencia de los cardinales de
  sus conjuntos potencia es $112$. Hallar $n(B)$.

  \medskip
  \noindent\textit{Planteamiento:}
  \[
    n[\mathcal{P}(B)] - n[\mathcal{P}(A)] = 112
    \;\Longrightarrow\;
    2^{n(B)} - 2^{n(A)} = 112
  \]
  \begin{align*}
    2^{7} - 2^{4} &= 128 - 16 = 112 \;\checkmark
  \end{align*}
  \begin{tcolorbox}[colback=AzulClaro!50, colframe=AzulPizarra,
      arc=2pt, boxrule=0.7pt]
    \[ \boxed{n(B) = 7} \]
  \end{tcolorbox}
\end{cajaejemplo}

\begin{cajaejemplo}[5 --- UNI: Conjunto universal y subconjuntos]
  Dado:
  \[
    U = \Bigl\{\tfrac{x+1}{3} \in \mathbb{N} \;\Big|\; 2 < x < 32\Bigr\},
    \quad P = \{x \in U \mid x \text{ es par}\}, \quad
    Q = \{x \in U \mid x < 7\}
  \]
  Calcular $n(P) - n(Q)$.

  \noindent\textit{Desarrollo:}
  \begin{align*}
    2 < x < 32 &\Rightarrow \tfrac{3}{3} < \tfrac{x+1}{3} < \tfrac{33}{3}
                \Rightarrow 1 < \tfrac{x+1}{3} < 11 \\
    U &= \{2, 3, 4, 5, 6, 7, 8, 9, 10\} \quad (n(U)=9)\\
    P &= \{2, 4, 6, 8, 10\} \;\Rightarrow\; n(P) = 5 \\
    Q &= \{2, 3, 4, 5, 6\}  \;\Rightarrow\; n(Q) = 5
  \end{align*}
  \begin{tcolorbox}[colback=AzulClaro!50, colframe=AzulPizarra,
      arc=2pt, boxrule=0.7pt]
    \[ \boxed{n(P) - n(Q) = 0} \]
  \end{tcolorbox}
\end{cajaejemplo}

% --- RESUMEN ---------------------------------------------------------------
\section{Resumen del tema}

\begin{cajaformula}
  \begin{center}
    \renewcommand{\arraystretch}{1.4}
    \begin{tabular}{@{}ll@{}}
      \toprule
      \rowcolor{AzulClaro!60}
      \textbf{Concepto} & \textbf{Expresi\'on / Descripci\'on} \\
      \midrule
      Cardinal         & $n(A) = $ n\'umero de elementos distintos \\
      Inclusi\'on      & $A \subset B \Leftrightarrow \forall\,x: x\in A \Rightarrow x\in B$ \\
      Conjunto potencia& $n[\mathcal{P}(A)] = 2^{n(A)}$ \\
      Uni\'on          & $n(A \cup B) = n(A)+n(B)-n(A\cap B)$ \\
      Diferencia sim.  & $n(A\mathbin{\triangle}B) = n(A)+n(B)-2n(A\cap B)$ \\
      Complemento      & $n(A') = n(U) - n(A)$ \\
      Vac\'{\i}o       & $\emptyset \subset A$, para todo conjunto $A$ \\
      \bottomrule
    \end{tabular}
  \end{center}
\end{cajaformula}

\begin{cajatip}
  \textbf{Estrategia clave para el examen:} cuando un problema mencione la
  ``diferencia de conjuntos potencia'', factoriza potencias de 2.
  Por ejemplo, $2^{n} - 2^{m} = 2^{m}(2^{n-m}-1)$. Esto acelera enormemente
  la identificaci\'on de $n(A)$ y $n(B)$.
\end{cajatip}

% FIN — teoria.tex
